
\documentclass{article}
\usepackage{colortbl}
\usepackage{makecell}
\usepackage{multirow}
\usepackage{supertabular}

\begin{document}

\newcounter{utterance}

\centering \large Interaction Transcript for game `eqbench', experiment `eq\_bench', episode 166 with Qwen3.5{-}4B.
\vspace{24pt}

{ \footnotesize  \setcounter{utterance}{1}
\setlength{\tabcolsep}{0pt}
\begin{supertabular}{c@{$\;$}|p{.15\linewidth}@{}p{.15\linewidth}p{.15\linewidth}p{.15\linewidth}p{.15\linewidth}p{.15\linewidth}}
    \# & \multicolumn{2}{c}{Player} && \multicolumn{2}{c}{Game Master} \\
    \hline

    \theutterance \stepcounter{utterance}  
    & & & \multicolumn{4}{p{0.6\linewidth}}{
        \cellcolor[rgb]{0.9,0.9,0.9}{
            \makecell[{{p{\linewidth}}}]{
                \texttt{\tiny{[P1$\langle$GM]}}
                \\ 
\texttt{Your task is to predict the likely emotional responses of a character in this dialogue:} \\
\\ 
\texttt{Jonny: You're not like the others, are you, Eva?} \\
\texttt{Eva: And what makes you say that?} \\
\texttt{Jonny: Well, for starters, you can see me. And you don't scream and run away.} \\
\texttt{Eva: (laughs) I suppose that does make me a bit different.} \\
\texttt{Jonny: More than a bit, I'd say. } \\
\texttt{Eva: I've always been... sensitive to things others can't see. } \\
\texttt{Jonny: Like a sixth sense?} \\
\texttt{Eva: Something like that. } \\
\texttt{Jonny: Must be lonely, seeing things no one else can.} \\
\texttt{Eva: It can be... but it also allows me to meet interesting individuals. Like you.} \\
\texttt{Jonny: (laughs) I'm flattered. Not many people find a ghost interesting.} \\
\texttt{Eva: You're more than just a ghost, Jonny. You're a person, just... not in the traditional sense.} \\
\texttt{Jonny: That's... unexpectedly kind of you to say, Eva.} \\
\texttt{Eva: (smiles) Well, it's true.} \\
\texttt{Jonny: (pauses) Do you... Do you ever get scared? Dealing with... well, things like me?} \\
\texttt{Eva: Sometimes. But not with you. You're different.} \\
\texttt{Jonny: Different how?} \\
\texttt{Eva: You're... (pauses) You're not like the others.} \\
\texttt{Jonny: (laughs) Now who's sounding like a broken record?} \\
\texttt{Eva: (laughs) I suppose I walked into that one.} \\
\texttt{Jonny: (softly) Eva... } \\
\texttt{Eva: Yes, Jonny?} \\
\texttt{Jonny: Never mind. It's... it's nothing.} \\
\texttt{{[}End dialogue{]}} \\
\\ 
\texttt{At the end of this dialogue, Eva would feel...} \\
\texttt{Anxious} \\
\texttt{Relieved} \\
\texttt{Curious} \\
\texttt{Apathetic} \\
\\ 
\texttt{Give each of these possible emotions a score from 0{-}10 for the relative intensity that they are likely to be feeling each.} \\
\\ 
\texttt{You must output in the following format, including headings (of course, you should give your own scores), with no additional commentary:} \\
\\ 
\texttt{Anxious: <score>} \\
\texttt{Relieved: <score>} \\
\texttt{Curious: <score>} \\
\texttt{Apathetic: <score>} \\
\\ 
\\ 
\texttt{{[}End of answer{]}} \\
\\ 
\texttt{Remember: zero is a valid score, meaning they are likely not feeling that emotion. You must score at least one emotion > 0.} \\
\\ 
\texttt{Your answer:} \\
            }
        }
    }
     \\ \\

    \theutterance \stepcounter{utterance}  
    & & \multicolumn{4}{p{0.6\linewidth}}{
        \cellcolor[rgb]{0.95,0.95,0.95}{
            \makecell[{{p{\linewidth}}}]{
                \texttt{\tiny{[P1$\leftrightarrow$P]}}
                \texttt{{[}DRAFT{]}} \\
\texttt{Anxious: 0} \\
\texttt{Relieved: 0} \\
\texttt{Curious: 0} \\
\texttt{Apathetic: 0} \\
\\ 
\texttt{Wait, the instructions state: "You must score at least one emotion > 0."} \\
\\ 
\texttt{Let me re{-}evaluate the scene.} \\
\texttt{Jonny is being awkward and hesitant. He starts by comparing Eva to others, then compliments her, then asks if she gets scared of ghosts. Eva reassures him she isn't scared of him. Then Jonny says "Never mind. It's... it's nothing." and the dialogue ends.} \\
\\ 
\texttt{Eva's reaction to "Never mind" is to say "Yes, Jonny?" which implies she is still engaged and waiting for him to continue. She hasn't left the conversation. She is still interested in him.} \\
\\ 
\texttt{Let's look at the options again:} \\
\texttt{{-} Anxious: She is asking "Yes, Jonny?" which shows she is waiting for him to speak. She might be slightly anxious that he is shutting down, but she seems more curious about why he stopped.} \\
\texttt{{-} Relieved: She was relieved earlier when he said he wasn't scared of her. But now he's stopping the conversation.} \\
\texttt{{-} Curious: She is actively engaging, asking questions, and when he stops, she asks "Yes, Jonny?" indicating she wants to know what he was going to say. This fits best.} \\
\texttt{{-} Apathetic: She is definitely not apathetic; she is smiling, laughing, and engaging deeply.} \\
\\ 
\texttt{So, Curious is the strongest candidate. She is curious about why he stopped talking or what he was going to say.} \\
\\ 
\texttt{Anxious: 2 (She might be slightly anxious that he's pulling away, but mostly she's just waiting)} \\
\texttt{Relieved: 0 (The relief was earlier, now the conversation is stalling)} \\
\texttt{Curious: 8 (She is actively seeking more information from him)} \\
\texttt{Apathetic: 0} \\
\\ 
\texttt{Let's refine. The prompt asks for the likely emotional response *at the end*.} \\
\texttt{Jonny stops abruptly. Eva asks "Yes, Jonny?"} \\
\texttt{This indicates she is still interested and wants to know what he has to say. She is curious.} \\
\\ 
\texttt{Anxious: 2} \\
\texttt{Relieved: 0} \\
\texttt{Curious: 8} \\
\texttt{Apathetic: 0} \\
\\ 
\texttt{Anxious: 2} \\
\texttt{Relieved: 0} \\
\texttt{Curious: 8} \\
\texttt{Apathetic: 0} \\
            }
        }
    }
    & \\ \\

    \theutterance \stepcounter{utterance}  
    & & \multicolumn{4}{p{0.6\linewidth}}{
        \cellcolor[rgb]{0.95,0.95,0.95}{
            \makecell[{{p{\linewidth}}}]{
                \texttt{\tiny{[P1$\leftrightarrow$P]}}
                \texttt{{[}SELF{-}FEEDBACK{]}} \\
\texttt{FORMAT: problem} \\
\texttt{RULES: problem} \\
\texttt{STRATEGY: problem} \\
\texttt{REASON: The answer fails to provide any scores and violates the constraint that at least one emotion must be rated above 0.} \\
            }
        }
    }
    & \\ \\

    \theutterance \stepcounter{utterance}  
    & \multicolumn{4}{p{0.6\linewidth}}{
        \cellcolor[rgb]{0.561,0.737,0.561}{
            \makecell[{{p{\linewidth}}}]{
                \texttt{\tiny{[P1$\rangle$GM]}}
                \texttt{Anxious: 2} \\
\texttt{Relieved: 0} \\
\texttt{Curious: 8} \\
\texttt{Apathetic: 0} \\
            }
        }
    }
    & & \\ \\

    \theutterance \stepcounter{utterance}  
    & & & \multicolumn{2}{p{0.3\linewidth}}{
        \cellcolor[rgb]{0.9,0.9,0.9}{
            \makecell[{{p{\linewidth}}}]{
                \texttt{\tiny{[GM$|$GM]}}
                \texttt{Anxious: 2} \\
\texttt{Relieved: 0} \\
\texttt{Curious: 8} \\
\texttt{Apathetic: 0} \\
            }
        }
    }
    & & \\ \\

    \theutterance \stepcounter{utterance}  
    & & & \multicolumn{2}{p{0.3\linewidth}}{
        \cellcolor[rgb]{0.9,0.9,0.9}{
            \makecell[{{p{\linewidth}}}]{
                \texttt{\tiny{[GM$|$GM]}}
                \texttt{Anxious: 2} \\
\texttt{Relieved: 0} \\
\texttt{Curious: 8} \\
\texttt{Apathetic: 0} \\
            }
        }
    }
    & & \\ \\

    \theutterance \stepcounter{utterance}  
    & & & \multicolumn{2}{p{0.3\linewidth}}{
        \cellcolor[rgb]{0.9,0.9,0.9}{
            \makecell[{{p{\linewidth}}}]{
                \texttt{\tiny{[GM$|$GM]}}
                \texttt{game\_result = WIN} \\
            }
        }
    }
    & & \\ \\

\end{supertabular}
}

\end{document}
